% 本文件由 LaTeX 排版 Agent 根据有效 translation.json 自动生成。
% author=Miscellaneous; work_id=0714; title=训诲

\chapter{\OriginalTerm{训诲}{Didache}}

% structure: p0001 source=h1 target=omitted reason=page-title-already-rendered-by-parent

\OriginalTerm{主的训诲}{Lord's Teaching}通过十二宗徒传给万民。\par

% structure: p0003 source=h2 target=section reason=relative-heading-rank; base=h2

\section{\OriginalTerm{两条道路}{Two Ways}；\OriginalTerm{第一条诫命}{First Commandment}}

有两条道路：一条是生命之道，一条是死亡之道；但两条道路之间有很大的差别。生命之道是这样：首先，你当爱那造你的天主；其次，爱邻人如同自己；凡你不愿发生在自己身上的事，也不要对别人做。这些教训的教导是：祝福那咒骂你们的，为你们的仇敌祈祷，为那迫害你们的禁食。因为若你们只爱那爱你们的人，有什么赏报呢？外邦人不也是这样行吗？但你们要爱那恨你们的，这样你们就没有仇敌。要戒绝肉体和世俗的情欲。若有人打你的右脸，把另一边也转给他，这样你就成为成全的。若有人强逼你走一里路，你就同他走两里。若有人拿走你的外衣，连内衣也给他。若有人拿走属于你的东西，不要追讨，因为你实在不能追讨。凡求你的，就给他，不要追讨；因为圣父愿意我们把自己的恩赐（白白赠予）分施给众人。按诫命施与的人是有福的，因为他无罪。领受的人有祸了；因为若那有需要的人领受，他是无罪的；但那没有需要却领受的人，要付出代价——他为何领受、为着什么，并且要陷入困境（拘禁），为自己所行的事受审，直到还清最后一分钱。\BibleRef{玛 5:26}但关于这事也说过：让你的施舍在你手中流汗，直到你知道该给谁。\par

% structure: p0005 source=h2 target=section reason=relative-heading-rank; base=h2

\section{\OriginalTerm{第二条诫命}{Second Commandment}：\OriginalTerm{粗重罪恶被禁止}{Gross Sin Forbidden}}

教训的第二条诫命：不可杀人，不可奸淫，\BibleRef{出 20:13}不可行男色，不可行淫乱，不可偷盗，\BibleRef{出 20:15}不可行邪术，不可行巫术，\OriginalTerm{堕胎}{abortion}不可杀害胎儿，不可杀害已生的婴孩。不可贪图你邻人的财物，\BibleRef{出 20:17}不可背誓，\BibleRef{玛 5:34}不可作假见证，\BibleRef{出 20:16}不可说恶言，不可怀怨。不可心怀二意，也不可一口两舌；因为一口两舌是死亡的网罗。你的言语不可虚假，也不可空洞，而要以行动成就。不可贪婪，不可强夺，不可伪善，不可怀恶意，不可傲慢。不可图谋恶事反对你的邻人。不可恨任何人；但要责备某些人，为某些人祈祷，并且爱某些人胜过自己的性命。\par

% structure: p0007 source=h2 target=section reason=relative-heading-rank; base=h2

\section{\OriginalTerm{其他被禁止的罪}{Other Sins Forbidden}}

我的孩子，你要逃避一切邪恶的事，以及一切类似邪恶的事。不可轻易发怒，因为愤怒导向杀人；也不可嫉妒、争吵、暴躁；因为这一切都滋生凶杀。我的孩子，不可成为好色之徒；因为色欲导向淫乱；也不可污言秽语、眼目高傲；因为这一切都滋生奸淫。我的孩子，不可观察兆头，因为这导向偶像崇拜；也不可行邪术、占星、洁净礼，也不可喜欢这些事；因为这一切都滋生偶像崇拜。我的孩子，不可说谎，因为谎言导向偷窃；也不可贪财、虚荣；因为这一切都滋生偷窃。我的孩子，不可抱怨，因为抱怨导向亵渎；也不可固执、恶意；因为这一切都滋生亵渎。但要温良，因为温良的人必承受土地。\BibleRef{玛 5:5}要忍耐、仁慈、纯朴、温和、良善，并始终敬畏你所听见的话语。不可自高，\BibleRef{路 18:14}也不可使你的灵魂过分自信。你的灵魂不可与高傲的人结交，却要与正义和谦卑的人交往。临到你的一切遭遇，都要当作好事领受，因为知道没有一事不是出于天主。\par

% structure: p0009 source=h2 target=section reason=relative-heading-rank; base=h2

\section{\OriginalTerm{各种训导}{Various Precepts}}

我的孩子，你要昼夜记念那向你讲论天主之道的人；并要尊敬他如同主，因为凡发出主言的地方，就有主。你要日日寻求圣人们的容面，以便安息在他们的言语中。不可渴求分裂，却要使争竞的人和好。你要秉公审判，不可在责备过犯时偏袒人。不可对事情是否要发生犹豫不决。不可伸手领取，却缩手不给。若你拥有什么，要藉你的手施舍作为你罪恶的赎价。不可迟疑给与，也不可抱怨；因为你知道谁是那善付工价者。不可转脸不顾那有缺乏的人，却要与你的弟兄共享一切，不可说东西是你的；因为你们既在不朽的事上有分，何况在必朽的事上呢？不可从你的儿子或女儿身上撤回你的手，却要从幼年起就教导他们敬畏天主。\BibleRef{弗 6:4}不可在你的怒气中命令你的男仆或婢女，因为他们指望同一位天主，恐怕你们所敬畏的天主不眷顾他们；\BibleRef{弗 6:9}因为祂来不是按外貌召叫，而是召叫那些圣神所预备的人。你们作仆人的，要怀着谦恭和敬畏，顺服你们的主人如同天主的样式。\BibleRef{弗 6:5}你要憎恨一切伪善，以及一切不悦乐主的事。万不可离弃主的诫命；却要持守你所领受的，既不增加，也不减少。\BibleRef{申 12:32}在教会中，你要承认你的过犯，不可带着恶的良心走近祈祷。这就是生命之道。\par

% structure: p0011 source=h2 target=section reason=relative-heading-rank; base=h2

\section{\OriginalTerm{死亡之道}{The Way of Death}}

死亡之路是这样：首先，它邪恶且充满咒骂：凶杀、奸淫、情欲、淫乱、偷窃、邪神崇拜、巫术、邪术、抢夺、假见证、伪善、心怀二意、欺骗、傲慢、堕落、自以为是、贪婪、污言秽语、嫉妒、自恃、高傲、夸口；迫害善人者、憎恨真理、喜爱谎言、不知义德的赏报、不依附善也不依附公义的判断、不为善事警醒，反为邪恶留意；远离温良和忍耐，喜爱虚荣，追求报复，不怜悯穷人，不为受苦者辛劳，不认识那造他们的主，杀害子女者，摧毁天主手工者，远离那有需要者，折磨那困苦者，富人的帮凶，穷人的不法审判官，彻底的罪人。孩子们，远离这一切吧。\par

% structure: p0013 source=h2 target=section reason=relative-heading-rank; base=h2

\section{第六章 反对假教师及祭偶像之物}

务要谨慎，免得有人使你们偏离这训言之道，因为这是出于天主的教训。因为如果你们能承受主全部的重轭，你们将会是成全的；如果不能，就尽力而为吧。至于食物，要量力而行；但对于祭过偶像的，要格外警惕，因为那是侍奉已死之神的事。\par

% structure: p0015 source=h2 target=section reason=relative-heading-rank; base=h2

\section{第七章 论圣洗}

此外，关于圣洗，要这样施行：首先宣讲这一切之后，因父、及子、及圣神之名给他们付洗\BibleRef{玛 28:19}，用活水。但如果没有活水，就用别的冷水；如不能用冷水，就用温水。如果两者都没有，就三次将水倒在头上，因父、及子、及圣神之名。但是在圣洗之前，施洗者应当禁食，受洗者也一样，以及其他可能的人；但要命令受洗者在之前禁食一日或两日。\par

% structure: p0017 source=h2 target=section reason=relative-heading-rank; base=h2

\section{第八章 论禁食与祈祷（主祷文）}

但你们的禁食不要与伪善者一样\BibleRef{玛 6:16}；因为他们在每周的第二日和第五日禁食，但你们要在第四日和预备日（星期五）禁食。也不要像伪善者那样祈祷，而要按主在他的福音里所命令的那样祈祷，这样祈祷：\BibleQuote{我们在天的父！愿你的名被尊为圣，愿你的国来临，愿你的旨意承行，如在天上，也在地上。求你今天赐给我们我们日用的食粮，并宽免我们的债，如同我们也宽免欠我们债的人；不要让我们陷入诱惑，但救我们脱离那恶者（或：邪恶）；因为权能和荣耀都是你的，直到永远。}这样每日祈祷三次。\par

% structure: p0019 source=h2 target=section reason=relative-heading-rank; base=h2

\section{第九章 感恩礼（圣体圣事）}

现在关于感恩礼（圣体圣事），要这样感恩。首先，关于杯：我们感谢你，我们的父，因你仆人\Person{达味}的圣葡萄树，你藉你的仆人\Person{耶稣}使我们认识它；愿光荣归于你，直到永远。关于擘开的饼：我们感谢你，我们的父，因生命和知识，你藉你的仆人\Person{耶稣}使我们认识它们；愿光荣归于你，直到永远。正如这擘开的饼曾散在山上，又聚合起来成为一体，同样，愿你的教会从地极聚集到你国里；因为光荣和权能藉\Person{耶稣基督}归于你，直到永远。但不要让任何人吃喝你们的感恩礼（圣体圣事），除非是那些因主的名受洗的人；因为主也对此说过：\BibleQuote{不要把圣物给狗}\BibleRef{玛 7:6}。\par

% structure: p0021 source=h2 target=section reason=relative-heading-rank; base=h2

\section{第十章 领圣体后的祈祷}

但在你们饱饫之后，要这样感恩：我们感谢你，圣父，因你的圣名，你使它住在我们心里，也因知识、信德和不朽，你藉你的仆人\Person{耶稣}使我们认识它们；愿光荣归于你，直到永远。全能的主宰，你为了你的名创造了万有；你赐给人饮食享受，好使他们感谢你；但给我们却藉你的仆人\Person{耶稣}白赐了属神的饮食和永生。在万有之前我们感谢你，因为你是大能的；愿光荣归于你，直到永远。上主，求你记念你的教会，救她脱离一切邪恶，并在你的爱中成全她；从四风把她聚集起来，圣化她，归入你所预备的国度；因为权能和光荣归于你，直到永远。愿恩宠来临，愿这世界逝去。达味之子（天主之子）和散那！凡圣洁的，让他前来；凡不圣洁的，让他悔改。主来临了。阿们。但让先知们随意感恩。\par

% structure: p0023 source=h2 target=section reason=relative-heading-rank; base=h2

\section{第十一章 论教师、宗徒及先知}

因此，凡前来教导你们上述所有事情的人，你们要接待他。但若教师本人转而教导另一种道理，以致于破坏这道理，你们不要听他；但如果他所教导的是增进义德和主的知识，就要接待他如同接待主一样。至于宗徒和先知，要按福音的规定这样行。凡来到你们那里的宗徒，要接待他如同主。但他不可停留超过一天；若有需要，也可第二天；但若他停留三天，他就是假先知。宗徒离开时，除了到达宿地所需的饼以外，不可拿任何东西；但若他索要金钱，他就是假先知。凡先知在圣神内讲话的，你们不可考验或判断；因为一切罪都将被赦免，但这一罪却不得赦免。但并非所有在圣神内讲话的都是先知，只有当他持守主的道路时才是。因此，从他们的道路就可以辨别假先知和真先知。凡先知在圣神内安排筵席的，他自己不可从中吃，除非他是假先知；凡先知教导真理，却不行他所教导的，他是假先知。凡先知经证明为真实，在世上为教会的奥迹劳苦，却不教导别人行他自己所行的，你们不可判断他，因为他在天主面前有自己的判断；因为古代的众先知也是如此行的。但若有人在圣神内说：给我钱或别的东西，你们不要听从他；但他若说为其他有需要的人而给，就不可判断他。\par

% structure: p0025 source=h2 target=section reason=relative-heading-rank; base=h2

\section{第十二章 接待基督徒}

但凡奉主名而来的人，你们要接待；其后你们要考验并认识他，因为你们将有左右分辨的智慧。若来者是旅客，要尽力帮助他；但他不可在你们那里停留超过两三天，除非有需要。若他愿意住在你们那里，且是工匠，就让他做工并吃饭；\BibleRef{得后 3:10}但若他没有手艺，按你们的判断，要留意他不可作为基督徒在你们中间闲居。若他不愿这样做，他就是贩卖基督者。要提防远离这样的人。\par

% structure: p0027 source=h2 target=section reason=relative-heading-rank; base=h2

\section{第十三章 供养先知}

但凡愿意住在你们中间的真先知，理应得到供养。同样，真正的教师，也如工人一样，理应得到他的供养。\BibleRef{玛 10:10}因此，凡酒榨和禾场上、牛群和羊群的初果，你们要取来给先知，因为他们是你们的大司祭。若你们没有先知，就应给穷人。若你和面，要取初果，按诫命奉献。同样，当你开一罐酒或油时，要取初果给先知；至于银钱、衣服和一切财物，也要取初果，按你们看为好的，并按诫命奉献。\par

% structure: p0029 source=h2 target=section reason=relative-heading-rank; base=h2

\section{第十四章 基督徒在主日的聚会}

每逢主日，你们要聚在一起，先告明你们的过犯，然后擘饼并感恩，好使你们的祭献得以纯洁。但凡与同伴不和的人，不可与你们一同聚集，直到他们和好，免得你们的祭献被亵渎。因为主曾说过：\Quote{在每一处所、每一时间，向我呈献纯洁的祭献；因为我是伟大的君王，——主说——我的名在万民中是奇妙的。}\par

% structure: p0031 source=h2 target=section reason=relative-heading-rank; base=h2

\section{第十五章 主教与执事；基督徒的规劝}

因此，你们要为自己委派配得上主的主教和执事，就是温良、不贪爱银钱\BibleRef{弟前 3:4}、且真诚、受过考验的人；因为他们也向你们履行先知和教师的职务。所以不可轻视他们，因为他们同先知和教师一样，是你们所尊敬的人。你们要彼此规劝，不可动怒，而要和平，正如你们在福音中所领受的；\BibleRef{玛 18:15}但凡亏负别人者，不要有人与他说话，也不可让他从你们那里听到任何话，直到他悔改。至于你们的祈祷、施舍和一切行为，都要照我们在主的福音中所领受的去行。\par

% structure: p0033 source=h2 target=section reason=relative-heading-rank; base=h2

\section{第十六章 警醒；主的来临}

要为你们的生命而警醒。不要让你们灯熄火灭，也不要让你们的腰束放松；但要时时准备好，因为你们不知道我们的主何时来临。\BibleRef{玛 24:42}你们要常常聚集，寻求那有益于你们灵魂的事；因为若你们不在末时得以成全，你们整个信德时期都将无益。因为在末日，假先知和败坏者将增多，羊将变成狼，爱将变成恨；\BibleRef{玛 24:11}因为不法之事增加，他们将彼此仇恨、迫害、出卖；\BibleRef{玛 24:10}然后那迷惑世界的假基督将显现，如同天主子，并施行奇事异能，大地要被交在他手中，他要做出自创世以来从未有过的不义之事。那时，人类将进入试炼之火，许多人将跌倒而丧亡；但那些在信德上坚忍的人，将得以从诅咒本身中获救。然后真理的记号将显现：首先是天上展开的记号，其次是号角的响声，第三是死者的复活；但不是所有的人，而是如经上所载：\Quote{主要带着他的众圣者同来。}那时，世界将看见主驾着天上的云彩降临。\par

\SourceRecord{Translated by M.B. Riddle.；Ante-Nicene Fathers；Vol. 7.；Edited by Alexander Roberts, James Donaldson, and A. Cleveland Coxe.；Buffalo, NY: Christian Literature Publishing Co.,；1886.；<http://www.newadvent.org/fathers/0714.htm>.}
